% 图 4.6 MNN 模型转换与动态推理 — 基于 CICC0903540 技术文档 4.4.1
\documentclass[tikz,border=8pt]{standalone}
% 顶会/顶刊风格：低饱和度马卡龙 + 高级感
% 适用于 NeurIPS, CVPR, ICLR, IEEE TPAMI
\usepackage{amssymb}
\usepackage{fontspec}
\usepackage{xeCJK}
\setCJKmainfont{SimHei}
\usepackage{tikz}
\usepackage{pgfplots}
\usetikzlibrary{
  arrows.meta, positioning, shapes.geometric, calc,
  backgrounds, fit, decorations.pathreplacing, patterns
}

% 低饱和度马卡龙配色
\definecolor{mBlue}{HTML}{AEC6CF}
\definecolor{mPink}{HTML}{FFD1DC}
\definecolor{mGreen}{HTML}{B1D8B7}
\definecolor{mPurple}{HTML}{B39EB5}
\definecolor{mYellow}{HTML}{FDFD96}
\definecolor{mGray}{HTML}{E5E4E2}
\definecolor{mDark}{HTML}{4A4A4A}

% 全局 TikZ 样式
\tikzset{
  >=stealth,
  every node/.style={align=center, font=\sffamily},
  block/.style={
    rounded corners=3pt, draw=mDark, align=center,
    minimum height=0.8cm, inner sep=5pt},
  arr/.style={-{Stealth[length=4pt]}, semithick, draw=mDark},
  % 信息密度增强：张量维度标注
  ts/.style={font=\tiny\sffamily, color=mDark!80},
  % 数学操作符节点（⊕加 ⊗乘 ‖拼接）
  op/.style={circle, draw=mDark, minimum size=0.4cm, inner sep=0,
    font=\scriptsize\sffamily, fill=mGray!60},
  % 图例
  legendbox/.style={rectangle, draw=mDark!50, rounded corners=2pt,
    fill=mGray!15, inner sep=4pt, font=\tiny\sffamily},
}

\begin{document}
\begin{tikzpicture}[
  x=1.05cm, y=1.0cm,
  node distance=0.5cm and 0.7cm,
  block/.style={
    rectangle, draw=mDark, fill=mBlue!40,
    minimum width=2.1cm, minimum height=0.8cm,
    font=\small\sffamily, align=center, rounded corners=4pt,
    semithick, inner sep=4pt},
  tool/.style={
    rectangle, draw=mDark!80, fill=mGray!50,
    minimum height=0.5cm, font=\scriptsize\sffamily\color{mGray},
    align=center, rounded corners=2pt},
  arr/.style={-{Stealth[length=4pt]}, semithick, draw=mDark},
]
  \node[block] (pt) {\texttt{.pt}};
  \node[tool, below=0.12cm of pt] {dynamic\_axes};
  \node[block, right=1.0cm of pt] (onnx) {ONNX};
  \node[tool, below=0.12cm of onnx] {MNN Converter};
  \node[block, right=1.0cm of onnx] (mnn) {\texttt{.mnn}};
  \node[tool, below=0.12cm of mnn] {Interpreter};
  \node[block, right=1.0cm of mnn] (resize) {resizeTensor};
  \node[tool, below=0.12cm of resize] {new\_shape};
  \node[block, right=1.0cm of resize] (sess) {resizeSession};
  \node[tool, below=0.12cm of sess] {内存预分配};
  \node[block, right=1.0cm of sess] (run) {runSession};

  \draw[arr] (pt) -- (onnx);
  \draw[arr] (onnx) -- (mnn);
  \draw[arr] (mnn) -- (resize);
  \draw[arr] (resize) -- (sess);
  \draw[arr] (sess) -- (run);

  \node[above=0.2cm of pt, font=\scriptsize\sffamily, color=mGray] {模型转换};
  \node[above=0.2cm of resize, font=\scriptsize\sffamily, color=mGray] {动态推理};

  % 图例（预留边距，张量标注已融入流程）
  \node[legendbox, anchor=north east] at ([xshift=-5pt, yshift=-5pt]current bounding box.south east)
    [align=left, text width=2.6cm] {
    \textbf{图例}\\[1pt]
    PyTorch $\to$ ONNX $\to$ .mnn
  };
\end{tikzpicture}
\end{document}
